
%%% M and N
\section{$\mathbf M$ and $\mathbf N$}
\label{sec:6.3}

Note that $\Omega=\avg K.$, the average element of the kernel. As well, $\Omega=\beta_\lambda\rho(k_\lambda)$,
where $k_i$ runs through any row, so that each range is represented once.

Let $N_0=\avg KK^*.  = \avg EE^*.$, where $K$ denotes a random element of $\cK$, and $E$ denotes a random idempotent,
element of $\E(\cK)$. Define $N=J-N_0$. $N$ satisfies the equation
%
\begin{equation}  \label{eq:N}
 \sym{A}{N}+\sym{\Delta}{N}=N
\end{equation} 
%
$N_0$ and $J$ satisfy the above equation as well.

Note that $N_0=\alpha_\mu e_\mu e_\mu^*$ where the $e_i$ are (all of) the idempotents with
a given range, i.e., they can be chosen as the idempotents from a single column in the kernel.

Let $M_0=\avg K^*JK.$. Define $\tau=\tr M_0$ and  $M=M_0-\tau\omega$. $M$ satisfies the equation
\begin{equation}\label{eq:M}
\symm{A}{M}+\symm{\Delta}{M}=M
\end{equation} 
$M_0$ and $\omega$ satisfy the above equation as well. Note that $JM_0=n^2\Omega$ and $\tr JM_0=n^2$.

Note that $M_0$, $N_0$, $M$, $N$ as well as $J$ are all symmetric matrices.
